Khi xảy ra sự cố rò rỉ khí trong nhà máy, kho chứa hay không gian hạn chế, hệ thống cảm biến cố định trả lời tốt câu hỏi \textit{``có rò rỉ hay không''} nhưng gần như không trả lời được câu hỏi \textit{``rò rỉ ở chỗ nào''}. Câu hỏi thứ hai hiện vẫn chủ yếu do con người giải quyết bằng cách trực tiếp mang máy dò đi vào vùng nguy hiểm.

Đề tài \textbf{GasSeeker} xây dựng một robot tự hành mang \textbf{một cảm biến khí duy nhất} để định vị nguồn phát khí trong một sân thử nghiệm 2 $\times$ 2 m, kèm một trạm mặt đất thu dữ liệu qua LoRa để người vận hành có thể đứng ngoài vùng nguy hiểm. Ràng buộc một cảm biến không phải là biện pháp cắt giảm chi phí mà là \textbf{điều kiện tạo ra bài toán}: khi không thể đo chênh lệch nồng độ trong không gian tại cùng một thời điểm, robot buộc phải dùng thời gian, chuyển động và thông tin hướng gió để thay thế. Báo cáo trình bày lập luận vật lý cho ràng buộc này, dựa trên cấu trúc đứt quãng của luồng khí trong dòng rối đã được công bố trong tài liệu và trên một phép phân tích tỉ số tín hiệu trên nhiễu cho thấy việc tăng khoảng cách giữa hai cảm biến trên cùng một xe không cải thiện được khả năng phân biệt hướng.

Nhóm cài đặt \textbf{ba thuật toán tìm nguồn} --- quét toàn bộ theo lưới, bám gradient nồng độ, và surge-casting mô phỏng hành vi bướm đêm --- trên \textbf{cùng một nền tảng phần cứng, cùng một quy trình đo và cùng một tiêu chí chấm điểm}, rồi so sánh chúng trong hai chế độ dòng khí: khuếch tán thuần và phát tán đứt quãng có gió. Toàn bộ mã thuật toán được tách thành một lớp C++ độc lập phần cứng, nhờ đó \textbf{chính đoạn mã sẽ nạp vào vi điều khiển cũng là đoạn mã chạy trong mô phỏng}.

Với 30 lần chạy mô phỏng cho mỗi tổ hợp thuật toán $\times$ môi trường (tổng cộng 180 lần chạy), kết quả cho thấy một hiệu ứng đảo chiều rõ rệt và có ý nghĩa thống kê: bám gradient đạt tỉ lệ thành công 67 \% trong môi trường khuếch tán nhưng chỉ còn 17 \% khi có gió (\textit{p} = 0,0002); surge-casting đi theo chiều ngược lại, từ 17 \% lên 57 \% (\textit{p} = 0,0028); trong khi quét toàn bộ giữ nguyên mức 77 \% và 63 \% (\textit{p} = 0,40) nhưng tốn gấp khoảng 2,2 đến 2,7 lần thời gian của hai thuật toán còn lại. Kết luận của đề tài \textbf{không phải} là ``thuật toán nào tốt nhất'' mà là một \textbf{quy tắc chọn thuật toán theo điều kiện dòng khí}.

Nhóm ghi nhận rõ giới hạn quan trọng nhất của báo cáo: \textbf{toàn bộ số liệu so sánh là kết quả mô phỏng, chưa phải số liệu đo trên xe thật}. Phần cứng đã được thiết kế, đấu nối và lập trình hoàn chỉnh, firmware biên dịch sạch, nhưng chưa kịp chạy đủ số lần thực nghiệm để đối chiếu.

\needspace{6\baselineskip}\noindent \textbf{Từ khoá:} định vị nguồn khí, robot khứu giác, cảm biến MQ-3, surge-casting, bám gradient, luồng khí rối, ESP32-S3, LoRa.
